\scanpage{15}

ramené au cas où $K$ correspond à une $p$-Algèbre de Lie nilpotente, puis, par le principe $2$, au cas où $K$ est commutatif, on gagne encore

b) $\dim K>0$. Par le principe $\unsure{2|1}$, \deleted{\unsure{donc|on}} on est ramené (par les commutateurs) au cas où $K$ est soit commutatif [qui est OK], soit $K=[K,K]$, qui implique $K$ affine. \deleted{\unsure{Enfin,}} Par le même principe [dévissant par un noyau de Frobenius, et utilisant le cas a)] on est ramené au cas où $K$ est simple$/k$.

\deleted{Enfin, \unsure{si on veut}, on peut le \unsure{supposer}}

Prenons alors, la section $\unsure{\varphi|\phi}$ de $\unsure{\sigma|\mathscr{T}}(K)$ correspondant aux tores maximaux, on \deleted{\unsure{donc}} se ramène \inserted{par le principe $1$} au cas où le tore maximal est \deleted{\unsure{donc}} invariant, donc par le principe $2$ aux \deleted{deux} \deleted{\inserted{trois}} \inserted{deux} cas suivants : a) $K$ fini ét., déjà traité b) $K$ \deleted{\unsure{comm.}} unip.\deleted{\unsure{otent}} connexe.

\deleted{Enfin, on peut dévisser par le \unsure{ss. groupe|ss-groupe} \unsure{radiciel|radical}, $=$ tore \unsure{maximal}, et on est \unsure{ramené}.} Alors on est ramené par dévissage par commutateurs, au cas comm., qui est OK.

% archive sheet number: 14

